
\documentclass{article}
\usepackage{colortbl}
\usepackage{makecell}
\usepackage{multirow}
\usepackage{supertabular}

\begin{document}

\newcounter{utterance}

\centering \large Interaction Transcript for game `textmapworld', experiment `small', episode 2 with reka{-}flash{-}3.1{-}t0.0.
\vspace{24pt}

{ \footnotesize  \setcounter{utterance}{1}
\setlength{\tabcolsep}{0pt}
\begin{supertabular}{c@{$\;$}|p{.15\linewidth}@{}p{.15\linewidth}p{.15\linewidth}p{.15\linewidth}p{.15\linewidth}p{.15\linewidth}}
    \# & $\;$A & \multicolumn{4}{c}{Game Master} & $\;\:$B\\
    \hline

    \theutterance \stepcounter{utterance}  
    & & \multicolumn{4}{p{0.6\linewidth}}{
        \cellcolor[rgb]{0.9,0.9,0.9}{
            \makecell[{{p{\linewidth}}}]{
                \texttt{\tiny{[P1$\langle$GM]}}
                \texttt{Please help me with the following task. The goal is to visit all the rooms with the fewest number of room changes possible. In each room, you need to decide the direction to go in. Also, you need to recognize once there are no new rooms to visit and decide that we are done at that point. Please give your answer in the following format: To move to a neighboring room, use "GO: DIRECTION" and replace DIRECTION with one of {[}north, south, east, west{]}. To stop the exploration, answer with "DONE" instead. Omit any other text.} \\
\texttt{Here is an example:} \\
\texttt{You are in the Kitchen. Currently available directions: south, west. What is your next instruction?} \\
\texttt{GO: west} \\
\texttt{You have made a step and entered a Lobby. Currently available directions: east, north. What is your next instruction?} \\
\texttt{GO: north} \\
\texttt{...} \\
\texttt{You have made a step and entered a Bedroom. Currently available directions: south. What is your next instruction?} \\
\texttt{DONE} \\
\texttt{Let us start. You are in the Pantry. Currently available directions: north. What is your next instruction?} \\
            }
        }
    }
    & \\ \\

    \theutterance \stepcounter{utterance}  
    & \multicolumn{4}{p{0.6\linewidth}}{
        \cellcolor[rgb]{0.9,0.9,0.9}{
            \makecell[{{p{\linewidth}}}]{
                \texttt{\tiny{[P1$\rangle$GM]}}
                \texttt{<reasoning>} \\
\texttt{Let's tackle this problem. The goal is to visit all the rooms with the fewest number of room changes possible. So, I need to figure out the optimal path that covers all rooms in the least steps. But how do I know the layout of the rooms? The user hasn't provided a map, so I must be working with the information given step by step as I move through the rooms.} \\
\\ 
\texttt{First, the starting point is the Pantry, and the only available direction is north. So the first move has to be north. Let's say I go north from the Pantry. Then I'll enter whatever room is there. Let's assume that's Room A. Now, I need to explore all connected rooms, but I have to track which rooms I've already visited to avoid cycles. Wait, but the problem says "visit all the rooms," so I need to make sure I don't get stuck in a loop and miss any rooms.} \\
\\ 
\texttt{But how do I know how many rooms there are in total? Without a map, it's tricky. Maybe the process is to explore each room, mark it as visited, and then proceed to unvisited rooms. But since the user is providing the available directions each time, I need to make decisions based on the current room's available exits.} \\
\\ 
\texttt{Wait, the example given shows that each time you move, you enter a new room, and the available directions are listed. So the assistant has to choose directions based on the current room's options. The challenge is to find the shortest path that covers all rooms. This sounds like a graph traversal problem, specifically a depth{-}first search (DFS) or breadth{-}first search (BFS). Since the goal is the fewest room changes (i.e., steps), BFS might be better because it explores all rooms at the current depth before moving to the next, which could lead to the shortest path. But DFS might get stuck in a long path if there's a dead end. However, without knowing the entire map, it's hard to apply BFS directly.} \\
\\ 
\texttt{Alternatively, maybe the optimal strategy is to always go to an unvisited room first, and if all adjacent rooms are visited, backtrack. This is similar to the "greedy" approach for maze solving. So, perhaps using a backtracking algorithm where you keep track of visited rooms and try unvisited directions first.} \\
\\ 
\texttt{But how to implement this without knowing the map? The assistant has to make decisions step by step, keeping} \\
            }
        }
    }
    & & \\ \\

    \theutterance \stepcounter{utterance}  
    & & & \multicolumn{2}{p{0.3\linewidth}}{
        \cellcolor[rgb]{0.9,0.9,0.9}{
            \makecell[{{p{\linewidth}}}]{
                \texttt{\tiny{[GM$|$GM]}}
                \texttt{abort game} \\
            }
        }
    }
    & & \\ \\

\end{supertabular}
}

\end{document}
